\section{Further research questions and the next structural targets}
\label{sec:research}

The current result reduces a coefficient, while the request for a smaller
order remains unresolved. The analysis gives sharper criteria for deciding
which next steps could actually change that outcome.

\paragraph{Actual image ranks of the crossing pieces.}
The four-energy obstruction rejects thinning the graded space while keeping
the same interfaces. It does not reject smaller factorizations through the
\emph{actual images} of the operators, provided a different, uniformly
proved factorization keeps the crossing identity intact. For a candidate
family $A_i$ on $\lambda$-bit strings, the target is
\[
 \max_i\left(\log_2\operatorname{rank}_{\rm row}A_i+
             \log_2\operatorname{rank}_{\rm col}A_i\right)
 \le4\rho\lambda+O(\log\lambda),\qquad \rho<1/\sqrt2,
\]
with polynomially many pieces, explicit integer factorizations,
polynomial-bit entries, and $2^{O(\lambda)}$ evaluation. An exact finite
rank test can reject bad candidates but cannot certify an asymptotic family.
The difficult new result would be a uniform cancellation or image-space
theorem, not a new Hilbert series substituted into the old proof.

\paragraph{More than two global regrouping directions.}
The Anthropic recurrence squares the clique size when it places the nodes
on a two-dimensional grid. A crossing identity compatible with three or
more useful global regroupings might change that recurrence. It would need
an explicit multi-direction identity, a balancing theorem, and counted
local evaluation. The two-order Heisenberg identity does not automatically
supply these. This is a structural research target, not a claimed algorithm.

\paragraph{Avoid explicit enumeration of correlated indices.}
Both the Cauchy mode sum and a hypothetical larger crossing-piece family
would benefit from structured summation of their indices. Merely encoding
an index by a short binary string does not give low-degree factors or a
cheap entry oracle. A proposed replacement must exhibit the factorization,
pay every coefficient query, retain creation precision, and prove a
strictly decreasing measure. The nested-clique calculation above explains
why composing the present two entropy estimates does not suffice.

\paragraph{Exploit correlations beyond the Euclidean envelope.}
The all-tree obstruction rules out improving the flag coefficient by
arbitrary tree imbalance alone. Actual generated mode factors could still
obey algebraic constraints absent from the relaxed entropy ball. A useful
lemma would show that the extremal star-heavy entropy profile cannot
arise together with costly child tensors, or would give an evaluator
sensitive to both remaining precision and mode weight. Simply lowering
requested precision leaves the inherited Hermite recovery increment and
does not by itself remove the square-root scale.

\paragraph{Improve certificates without confusing them with order gains.}
The pair certificate is a low-order relaxation of a finite support-choice
problem. Distributing the norm over triples or larger subsets, or choosing
another rational allocation of diagonal terms among pairs, can give tighter
certificates. The proof generalizes by counting how often each diagonal and
cross term appears. Such work could further lower explicit constants and
improve resource tradeoffs. As long as the mode entropy remains
$M/\sqrt u$ with $u=\Theta(\log s)$, it will not change the logarithmic order.

\paragraph{Formalize the small new lemmas first.}
The norm-distribution identity, finite support generation, all-tree recurrence
inequalities, and rank-window theorem are compact formalization targets.
They can be verified separately from the large analytic contractor. A later
end-to-end integration should then connect those lemmas to the exact ring,
packing, and machine-cost interfaces rather than treating a successful
finite checker as verification of the whole SAT algorithm.
